% Moved out of the main text for length only. Nothing was deleted;
% these sections are reported in supplementary material.

\subsection{A second domain, and what it took to build one}
\label{sec:synthfail}

We attempted a second oracle-bearing domain, a program updating integer variables
with interleaved queries, to test whether these findings are properties of
long-horizon state tracking or of chess. Four designs failed to leave chance, and
we report the sequence because the diagnosis is more useful than the intention.

The first placed a single query at the end of a sequence of random updates.
Training loss settled at the token-type entropy of the generator, which is the
value a model achieves by learning the format and nothing else, and the reason is
arithmetic: the only learnable signal was one answer token in roughly one
hundred and thirty, so state tracking carried under one percent of the gradient.
Interleaving queries raised that share and did not move accuracy. Reducing the
modulus and the number of variables did not move it either. A final version
separated a pure read-out query from a computed one, so that state tracking could
succeed without the arithmetic succeeding, and read-out accuracy still sat at
chance.

We draw two conclusions from the arithmetic attempts. The narrow one is that this task
family asks a small model to learn modular arithmetic before any state tracking
becomes observable, and that our compute did not reach it. The broader one is a
warning for anyone building synthetic state-tracking benchmarks: a task can be
perfectly well specified, carry an exact oracle, and still measure nothing,
because the quantity of interest contributes almost none of the training signal.
The loss curve diagnoses this cheaply. A model that has learned the format and no
content sits at the generator's token-type entropy, which is computable in
advance and worth computing before trusting such a benchmark. We report this in
place of an easy cross-domain result.

\paragraph{A domain that does work.}
Abandoning arithmetic entirely, we built the second domain around an assignment
instead. Four objects sit in four boxes. The sequence states the initial
assignment, then applies operations that permute it, and after every operation
asks where one object now is. A \textsc{swap} names two boxes and moves their
contents without naming the objects, so the affected object's location cannot be
retrieved by attending to its last mention and has to be maintained. A
\textsc{move} names the object, which gives an anchor and keeps the task
learnable. The state is the assignment of objects to boxes, exact at every
token, and the answer is a single symbol, so belief consistency here is exact
rather than estimated: a wrong answer either is or is not the box the model's own
probe says holds that object.

Two further design choices were forced on us, and both are worth recording. Our
first attempt used five objects. That makes the state an element of $S_5$, a
permutation-composition problem that is theoretically hard for the relevant
fixed-depth model classes under asymptotic expressivity assumptions
\citep{merrill2024illusion}. That may have contributed, but a theorem about
asymptotic expressivity does not establish why one finite model failed to
optimise on one bounded-length dataset, and we do not claim group structure as
the cause. Four objects give $S_4$,
which is solvable. Our second attempt queried only occasionally, which left the
task all-or-nothing: with uniformly random updates and sparse queries a model
either tracks the permutation exactly or sits at chance, a structure that is hard
to learn from endpoint supervision. Querying after every operation supplies
partial credit, since a query one step after an operation is answerable directly,
and that is the foothold from which the model extends. Chess has this property
for free, because a plausible move needs only partial state.

With those two changes the task became measurable, and the outcome follows.

The outcome is not the clean replication we hoped for, and it is more
informative than one would have been (Table~\ref{tab:boxes}).

With the supervised recurrent channel the task is close to solved at every depth,
\result{ansDeepBoxAlsb} at the deepest bucket against a chance rate of
\result{chanceBoxAlsb}, the probe recovers the queried object at
\result{probeDeepBoxAlsb}, and belief consistency is \result{beliefBoxAlsb}
against a mismatched control of \result{beliefMisBoxAlsb}. The mechanism
reproduces.

The attention-only model behaves differently. It solves the task well above
chance, \result{ansDeepBoxAttn} at depth, while the probe recovers its state at
\result{probeDeepBoxAttn}, which is chance. Belief consistency is
\result{beliefBoxAttn} against a control of \result{beliefMisBoxAttn}, which is
no signal at all. The same probe procedure, on the same positions, recovers state
at \result{probeDeepBoxAlsb} in the other model, so this is a fact about the
model and not a failure of the probe.

The reading we prefer is that this task does not force the model to maintain
state. Only one object is asked about at a time, so the answer can be resolved on
demand by tracing that object's history, and an explicit assignment carried
forward is unnecessary. Chess allows no such shortcut: predicting a legal move
requires the whole position, so the state must be maintained eagerly.

That gives a fourth lesson for anyone building these diagnostics, and it is the
one we would put first. A task can carry an exact oracle, be learnable, and still
fail to measure state maintenance, because it never obliges the model to
maintain anything. Belief consistency is only defined where a belief exists to
read, and whether one exists is a property of the task, not only of the
architecture.



\subsection{Is the measurement useful, or only diagnostic?}
\label{sec:abstain}

Everything so far is a claim about mechanism. A reader is entitled to ask whether
the same measurement is worth anything operationally, and the only baseline that
matters is the model's own confidence. If an external state monitor merely
recovers what the model already knows about its own uncertainty, it is
interesting and redundant.

We test this as selective prediction. A system that wants to avoid acting on a
decayed state abstains on the positions it trusts least, and the question is what
to rank by. We compare the model's top-1 probability, the entropy of its
next-move distribution, and the probe's certainty about the squares the chosen
move touches. That last signal uses only the probe's own output, so it is
computable at inference from activations, the probe having been fitted against an
oracle offline.

One version of this experiment is not admissible and we report it only as a
bound. Ranking by whether the probe is \emph{wrong} on those squares requires the
oracle at inference, which is the very thing an abstention rule exists to avoid.
It is an upper bound on what a perfect state monitor could deliver, not a method.

The honest result is modest. Probe certainty alone is a weaker signal than the
model's own confidence, with an AUC of \result{abAucProbe} against
\result{abAucConf}. Combined with confidence it is better than either alone: at
half coverage the illegal-move rate among retained positions falls to
\result{abComb50} against \result{abConf50} for confidence and
\result{abEnt50} for entropy, and in a joint fit both terms carry weight,
\result{abBetaConf} for confidence and \result{abBetaProbe} for probe certainty.
An external monitor therefore adds something the model does not self-report, and
it adds less than we hoped.

The gap to the bound is the interesting part. Ranking by oracle-verified
move-touched state reaches \result{abOracle50} at the same coverage, less than
half the base rate of \result{abBase}. Most of the benefit an ideal state monitor
would deliver is not captured by the probe's own certainty, which says the
limitation is the quality of the state read-out rather than the premise. Better
read-outs, not a different signal, are the lever.


\subsection{Transcript length against position difficulty}
\label{sec:cycles}

The decay curves are read throughout as loss caused by composing more updates.
Depth is confounded with position difficulty: a position at ply 100 differs from
one at ply 10 in material, phase, and legal-move count. The curves alone do not
license the causal reading, and this test removes the confound.

We insert a reversible cycle, a knight out and back for each side, which returns
the position exactly while adding four plies of transcript. State is then read on
an identical position at several transcript lengths.

\begin{table}[t]
\centering\small
\caption{Transcript length with the position held fixed. A reversible
cycle of legal moves returns the position exactly while adding four plies
of history, so every row scores the same states at a different transcript
length. Intervals are bootstrapped over games.}
\label{tab:cycles}
\begin{tabular}{lccc}
\toprule
Added plies & occupied-square acc. & 95\% CI & illegal rate \\
\midrule
0 & 0.736 & [0.722, 0.747] & 0.065 \\
4 & 0.715 & [0.702, 0.727] & 0.050 \\
8 & 0.706 & [0.692, 0.718] & 0.113 \\
12 & 0.696 & [0.683, 0.707] & 0.177 \\
\bottomrule
\end{tabular}
\end{table}


Fidelity falls monotonically, from \result{cycAcc0} [\result{cycLo0},
\result{cycHi0}] with no insertion to \result{cycAcc12} [\result{cycLo12},
\result{cycHi12}] after twelve added plies, on \result{cycN} positions whose
state is unchanged. The intervals at the two ends do not overlap. Transcript
length alone therefore causes measurable state loss.

The size matters as much as the sign. Twelve plies of pure length cost about four
points of accuracy, where the natural depth curve falls far further over its
range. Both factors contribute, and we do not attribute the whole curve to
composition depth. One caveat we cannot remove: the inserted moves are legal but
unlike training data, so part of the drop may be distribution shift rather than
length as such. The random-play result of Section~\ref{sec:randomplay} bounds
that concern without eliminating it.


\subsection{Inference that respects the clustering}
\label{sec:stats}

Positions within a game are not independent: they share an opening, a player, and
an accumulating history. Intervals computed over positions therefore understate
uncertainty, and every interval in this section resamples whole games instead.
Section~\ref{sec:belief} already reports game-clustered intervals for belief
consistency on the full evaluation population. This section is a restricted
sensitivity analysis on a different population, \result{stN} positions from
\result{stGames} games of which \result{stIll} carry an illegal top-1 move, over
ply 10 to 80 where the depth buckets used by the predictor comparison are
defined. Estimates therefore differ slightly from the primary ones, and both are
reported rather than reconciled by choosing whichever is larger.

\begin{table}[t]
\centering\small
\caption{Inference clustered on the game. Intervals resample whole
games rather than positions, since positions within a game share an
opening and a history. The lower block compares nested models by held-out
log loss on a split by game, so each predictor's contribution is measured
out of sample. Lower is better.}
\label{tab:stats}
\begin{tabular}{lcc}
\toprule
Quantity & Estimate & 95\% CI (clustered) \\
\midrule
belief consistency & 0.3892 & [0.3822, 0.3976] \\
\quad mismatched-belief control & 0.0654 & [0.0620, 0.0687] \\
\quad random-move control & 0.0079 & [0.0068, 0.0093] \\
AUC, whole board & 0.6287 & [0.6233, 0.6333] \\
AUC, move-touched & 0.6961 & [0.6926, 0.6998] \\
\quad difference & +0.0674 & [+0.0618, +0.0732] \\
\midrule
Held-out log loss & \multicolumn{2}{c}{} \\
depth only & 0.3630 & \\
\quad $+$ whole board & 0.3595 & \\
\quad $+$ move-touched & 0.3451 & \\
\quad $+$ both & 0.3439 & \\
\quad $+$ both and confidence & 0.3340 & \\
\bottomrule
\end{tabular}
\end{table}


Belief consistency is \result{cbBel} [\result{cbBelLo}, \result{cbBelHi}], against
\result{cbBelMis} [\result{cbBelMisLo}, \result{cbBelMisHi}] for a mismatched
belief and \result{cbBelRand} [\result{cbBelRandLo}, \result{cbBelRandHi}] for a
random illegal move. The intervals are far apart, and the separation does not
depend on treating positions as independent.

The two predictors of an illegal move separate as well. Whole-board error reaches
an AUC of \result{cbAucAgg} [\result{cbAucAggLo}, \result{cbAucAggHi}] and
move-touched error \result{cbAucTouch} [\result{cbAucTouchLo},
\result{cbAucTouchHi}]. Bootstrapping the difference as a paired quantity gives
\result{cbAucDiff} [\result{cbAucDiffLo}, \result{cbAucDiffHi}], which excludes
zero.

A population-averaged logistic fit, with game as the cluster and an exchangeable
working correlation, gives coefficients with clustered standard errors on standardised
predictors: \result{geeTouch} (clustered s.e.\ \result{geeTouchSE}) for move-touched
error against \result{geeAgg} (\result{geeAggSE}) for the aggregate, with
\result{geeConf} for low model confidence and \result{geeDepth} for depth. The
ordering survives clustered inference.

\paragraph{Out of sample, not only in sample.} Coefficients can be significant and
still add little. Splitting by game and fitting on seventy percent of games, the
held-out log loss falls from \result{lldepthonly} with depth alone to
\result{llplusaggregate} adding whole-board error, and to \result{llplustouched}
adding move-touched error instead. Using both gives \result{llplusboth}, and
adding the model's own confidence \result{llplusbothandconfidence}. The
move-touched predictor buys several times what the aggregate buys out of sample,
and it buys most of it even when the aggregate is already present.

\paragraph{A held-out abstention rule.} Fitting and evaluating a selective
prediction rule on the same positions overstates it. Refitting on training games
and evaluating on held-out games, the illegal rate among retained positions at
half coverage is \result{hoComb05} for a rule combining the model's confidence
with probe certainty, against \result{hoConf05} for confidence alone and
\result{hoCert05} for probe certainty alone, from a base rate of
\result{hoBase}. An oracle-verified state monitor would reach
\result{hoBound05}, which remains the bound rather than a method: establishing it
at inference needs the ground truth the rule exists to avoid.


\subsection{A policy-conditioned state-fidelity horizon}
\label{sec:hfid}

Both horizons used so far are read off behaviour, and the read-out ablation
showed why that is fragile: a model can lower its illegal-move rate by hedging
toward the legal-move prior, without knowing the position any better. A horizon
defined on the representation is less exposed, because it does not read the
confidence of the output: a model that spreads its mass more evenly gains nothing
from doing so. It is not policy-independent, and we do not claim it is. The
squares it scores are those the selected move touches, and that move comes from
the output distribution, so a policy change can alter which squares are scored
even with the state representation unchanged. It is a policy-conditioned
state-fidelity horizon. A genuinely policy-independent comparison would score a
fixed set of components, for instance the recorded continuation move, every
oracle-legal move averaged uniformly, or a candidate set shared across models.

We define $H_{\mathrm{fid}}$ as the depth at which the probability that the probe
is wrong on the squares the model's chosen move touches first crosses one half.
It uses the same probe and the same activations as everything else, and it is
immune to hedging: a model that spreads mass over legal moves gains nothing,
since the measure asks what it knows, not what it emits.

Read this way the ladder is orderly. The smallest model reaches
\result{hfidRung1} ply, the baseline \result{hfidAttn}, and the largest
\result{hfidRung6}. The behavioural horizon orders the same models the same way,
\result{hintRung1}, \result{hintAttn} and \result{hintRung6} ply, on a shorter
scale because it asks a harder question of the model. Move-touched error at
the shallowest depth runs from \result{actErrFirstRung1} in the smallest model to
\result{actErrFirstRung6} in the largest. We report both horizons throughout, and
prefer this one where the comparison is between architectures rather than between
depths, since that is where hedging is most likely to confound a behavioural
reading.


\section{Experimental detail}
\label{app:methods}

\paragraph{Data.} Public Lichess archives, standard chess only, both players rated
at least 1600, games between 30 and 160 ply, deduplicated by a hash of the move
sequence. Moves are tokenised one per move in UCI form, giving a vocabulary of
\result{nVocab} including padding and a start symbol; there is no sub-move
tokenisation and no board rendering in the input. Sequences are padded to 161
tokens and truncated at 160 ply. Splits are by game: \result{ntrainlm} for
language-model training, \result{ntrainprobe} for probe fitting,
\result{neval} held out for evaluation.

\paragraph{Oracle.} Every game is replayed with a rules engine, recording after
each token the occupancy of all 64 squares as one of 13 classes, the side to
move, and castling rights. Replay completeness is verified before any model is
measured.

\paragraph{Models.} Pre-norm decoder-only Transformers with learned positional
embeddings, GELU feed-forward of width four times the model dimension, and no
dropout. The grid spans 6 to 12 layers and widths 192 to 512, with 6 or 8 heads.
Training uses AdamW, $\beta = (0.9, 0.95)$, weight decay 0.01, a one-cycle
schedule peaking at $3\times10^{-4}$ with 5 percent warmup, batch size 32, 12000
steps, and mixed precision with gradient clipping at 1.0. The final checkpoint is
used; no checkpoint selection is performed. Seeds 0, 1 and 2 control
initialisation, batch order, and any auxiliary sampling.

\paragraph{Probes.} One linear head per layer maps the residual stream to 64
squares by 13 classes, trained with cross-entropy on the probe split only, AdamW
at $10^{-3}$, two epochs, batch size 48. The reported layer is the one with the
highest occupied-square accuracy, selected on the probe split rather than on the
evaluation split. Controls are a per-square per-depth majority predictor, a
label-permutation control pairing activations with another game's state at equal
depth, and a randomised-weight model of identical architecture whose probes are
fitted the same way.

\paragraph{Believed boards.} A decoded board is built by placing the argmax class
at each square, setting side to move from ply parity, and clearing castling
rights so that castling cannot be licensed by an assumption we did not decode.
Legality is tested as pseudo-legality, since a decoded board may lack a king and
make check tests ill-defined. Moves that fail to parse are counted as illegal.

\paragraph{Engine scoring.} Stockfish 17 at fixed depth 10. Loss is the drop from
the best available score to the score after the played move, both from the
mover's point of view, with mate scores mapped to 10000. The blunder threshold is
100 centipawns, fixed in advance.

\paragraph{Interventions.} Edits are added to the residual stream at the output of
the probed block, at the final position only. The direction is the difference of
probe decoder rows for the target and current classes, normalised and scaled to
$\alpha$ times the residual stream's own root-mean-square magnitude times the
square root of the model dimension. The reported $\alpha$ is \result{patchAlphaAttn},
chosen from the sweep in Table~\ref{tab:patchcal} as the smallest value that
reliably flips the decoded belief while leaving a norm-matched random edit at
chance. Random controls are drawn isotropically and rescaled to the same norm.

\paragraph{Recurrent channel.} A gated diagonal recurrence of width 64 inserted
after every fourth block, with the read-out matrix zero-initialised so the module
begins as the identity. The auxiliary objective, when active, is cross-entropy
from the channel to the oracle state with weight 1.0. The parameter-matched
control widens the feed-forward multiplier to 4.156 to absorb the same parameter
count.

\paragraph{Second domain.} Four objects in four boxes, an explicit initial
assignment, then 24 operations drawn evenly between a swap of two named boxes and
a move of a named object to a named box, with a query after every operation
asking where one object is. Answers are weighted five times in the loss. Details
of the six designs that failed before this one are in Section~\ref{sec:synthfail}.
